% GENERATED FILE. Edit canonical lessons, then run npm run generate:books.
% canonical-source-hash: fnv1a64:dd4fd259b7f35fa5
% canonical-lessons: SA-C59-desha, SA-C59-bharatam, SA-C59-bharatiyah, SA-C59-bharatat

\chapter{Where You Come From}
\label{ch:sa-where-you-come-from}
\hlchaptermodality{59}

\begin{chapteropening}
\textbf{By the end of this chapter:} \emph{I can name a country, name India, say that somebody is Indian, and answer \sk{कुतः} with the ablative.}

Complete SA-C59-bharatat and answer \sk{कुतः} out loud, then say that he comes from India.
\end{chapteropening}

\section[deśaḥ]{\sk{देशः} (deśaḥ) — a country}

\label{lesson:SA-C59-desha}



\begin{quote}
Say \textbf{\sk{कुतः}} and what it asks. Then say \textquotedblleft{}I like the fruit\textquotedblright{}. You have been able to ask \emph{from where?} since the question words and had nothing to answer it with.
\end{quote}

\subsection*{You'll want to know: \sk{देशः}}
\textbf{\sk{देशः}} (\emph{deśaḥ}) — \textquotedblleft{}a country, a region\textquotedblright{}.

The track has taught you a village, a city, a road and a forest. It has never taught you the largest of the containers, and it is the one a form asks for first.

The word is built on \textbf{\sk{दिश्}} (\emph{diś-}), \textquotedblleft{}to point out, to show\textquotedblright{}. A \textbf{\sk{देशः}} began as whatever a pointing hand took in — a stretch of land indicated rather than measured.

That root travelled a long way west. Latin \emph{dīcere}, \textquotedblleft{}to say\textquotedblright{}, and \emph{index}, the pointing finger, grow from the same ancient root as \textbf{\sk{दिश्}}. The finger you use to show somebody where you are from is the finger the word is named after.

\begin{grammarlens}[title={What you've built}]
The first of three words for saying where you are from.
\end{grammarlens}

\subsection*{Your turn}
\begin{itemize}
\raggedright
  \item \emph{Recall:} say \emph{grāmaḥ}, then say \textbf{\sk{कुतः}} and what it asks
  \item \emph{Hear:} \emph{deśaḥ}, said once slowly and once at speed
  \item \emph{Say it:} \emph{deśaḥ}
  \item \emph{Read it:} \textbf{\sk{देशः}}, then say what it means
\end{itemize}

\begin{tcolorbox}[breakable,colback=teal!4,colframe=teal!35!black,title={Before you move on}]
What does \textbf{\sk{देशः}} mean, and which root is it built on? (\textquotedblleft{}A country\textquotedblright{}, and \textbf{\sk{दिश्}}, to point out.)
\end{tcolorbox}
\section[Bhāratam]{\sk{भारतम्} (Bhāratam) — India}

\label{lesson:SA-C59-bharatam}



\begin{quote}
Say \emph{deśaḥ} and what it means. Then say \textbf{\sk{अस्ति}}.
\end{quote}

\subsection*{You'll want to know: \sk{भारतम्}}
\textbf{\sk{भारतम्}} (\emph{Bhāratam}) — \textquotedblleft{}India\textquotedblright{}.

A country and the name of one:

\begin{quote}
\textbf{\sk{भारतं} \sk{देशः} \sk{अस्ति}} — \emph{Bhārataṁ deśaḥ asti} — \textbf{India is a country}
\end{quote}

The name is older than any modern country. \textbf{\sk{भरत}} (\emph{Bharata}) is a figure the Sanskrit epics name, and \textbf{\sk{भारत}} is the form built from it meaning \textquotedblleft{}belonging to Bharata\textquotedblright{} — the land of his people. The great epic is called the \textbf{\sk{महाभारतम्}}, and the \textbf{\sk{महा}-} on the front of it is the \textbf{\sk{महत्}} you learnt in the adjectives: \emph{the great {[}story{]} of the Bharatas}.

English \emph{India} took an entirely different road. It came through Greek and Persian from \textbf{\sk{सिन्धु}} (\emph{Sindhu}), the river — the same word that gives \emph{Sindh} and, by a further turn, \emph{Hindu}. So one country carries two names that began as two different things: a people and a river.

\begin{grammarlens}[title={What you've built}]
A country, and the name of one.
\end{grammarlens}

\subsection*{Your turn}
\begin{itemize}
\raggedright
  \item \emph{Say it:} \emph{Bhāratam}
  \item \emph{Say it:} \emph{Bhārataṁ deśaḥ asti}
  \item \emph{Recall:} read \textbf{\sk{देशः}} and say what it means
  \item \emph{Read it:} \textbf{\sk{महाभारतम्}}, and find the \textbf{\sk{महा}-} you already know inside it
\end{itemize}

\begin{tcolorbox}[breakable,colback=teal!4,colframe=teal!35!black,title={Before you move on}]
What is \textbf{\sk{भारतम्}}, and where does the name come from? (India, from \textbf{\sk{भरत}}.) Say \textquotedblleft{}India is a country\textquotedblright{}.
\end{tcolorbox}
\section[bhāratīyaḥ]{\sk{भारतीयः} (bhāratīyaḥ) — Indian}

\label{lesson:SA-C59-bharatiyah}



\begin{quote}
Say \textbf{\sk{भारतम्}}. Then recall the adjective rule: where does an adjective stand, and what does it do to match its noun?
\end{quote}

\subsection*{You'll want to know: \sk{भारतीयः}}
\textbf{\sk{भारतीयः}} (\emph{bhāratīyaḥ}) — \textquotedblleft{}Indian\textquotedblright{}.

A country name is a noun. Saying that a \textbf{person} belongs to it needs an adjective, and Sanskrit makes one with a suffix you can now watch working:

\begin{quote}
\textbf{\sk{भारतम्}} $\to$ \textbf{\sk{भारतीयः}} — \emph{of Bhārata, Indian}
\end{quote}

\begin{quote}
\textbf{\sk{सः} \sk{भारतीयः}} — \emph{saḥ bhāratīyaḥ} — \textbf{he is Indian}
\end{quote}

The piece that does the work is the \emph{-īya} suffix, and you can see it in the word: the \textbf{\sk{ी}} after \textbf{\sk{भारत}}, then \textbf{\sk{य}}, then the masculine \textbf{-\sk{ः}}. It is a machine, not a one-off. Give it a place and it gives you the adjective for belonging to that place, and the ending then agrees with whoever you are describing exactly as \textbf{\sk{महत्}} and \textbf{\sk{नव}} do.

Notice what the sentence does \textbf{not} have. There is no \textbf{\sk{अस्ति}} in it. Sanskrit can leave the \emph{is} out when it simply links two things, and here it does — \emph{he, Indian}. You have seen that before in \textbf{\sk{न} \sk{चिन्ता}}; this is the same economy.

\begin{grammarlens}[title={What you've built}]
A country, its name, and the word for somebody who belongs to it.
\end{grammarlens}

\subsection*{Your turn}
\begin{itemize}
\raggedright
  \item \emph{Say it:} \emph{bhāratīyaḥ}
  \item \emph{Say it:} \emph{saḥ bhāratīyaḥ}
  \item \emph{Recall:} say \textbf{\sk{भारतम्}}, then build the adjective from it out loud
  \item \emph{Read it:} \textbf{\sk{भारतीयः}}, and find the \textbf{\sk{ी}} and the \textbf{\sk{य}} that build it
\end{itemize}

\begin{tcolorbox}[breakable,colback=teal!4,colframe=teal!35!black,title={Before you move on}]
What suffix turns a place into \textquotedblleft{}belonging to that place\textquotedblright{}? (\emph{-īya}.) Say \textquotedblleft{}he is Indian\textquotedblright{}.
\end{tcolorbox}
\section[Bhāratāt]{\sk{भारतात्} (Bhāratāt) — from India}

\label{lesson:SA-C59-bharatat}



\begin{quote}
Say \textbf{\sk{कुतः}}. Then say \textquotedblleft{}he is Indian\textquotedblright{}.
\end{quote}

\begin{grammarlens}[title={Grammar Lens: the ending that answers the question}]
\textbf{\sk{कुतः}} has been in this book since the question words asking \emph{from where?} and until now the book had no way to answer it. Naming the country was only half the fix. The other half is an ending.

\begin{quote}
\textbf{\sk{कुतः}?} — \emph{from where?}
\end{quote}

\begin{quote}
\textbf{\sk{भारतात्।}} — \emph{from India.}
\end{quote}

The \textbf{-\sk{आत्}} is the \textbf{ablative} — the case whose one job is \emph{away from}. Put it on a place and the place becomes a starting point:

\noindent
\begin{tabularx}{\linewidth}{@{}>{\raggedright\arraybackslash}X>{\raggedright\arraybackslash}X>{\raggedright\arraybackslash}X@{}}
\toprule
\textbf{word} & \textbf{ending} & \textbf{meaning} \\
\midrule
\textbf{\sk{भारतम्}} & -\sk{अम्} & India \\
\textbf{\sk{भारतात्}} & -\sk{आत्} & from India \\
\textbf{\sk{देशात्}} & -\sk{आत्} & from the country \\
\bottomrule
\end{tabularx}

And in a whole sentence, with the come-verb you already hold:

\begin{quote}
\textbf{\sk{सः} \sk{भारतात्} \sk{आगच्छति}}
\end{quote}

\begin{quote}
\emph{saḥ Bhāratāt āgacchati}
\end{quote}

\begin{quote}
\textbf{he comes from India}
\end{quote}

This is the pattern to keep, because it is how Sanskrit answers most questions of place: \textbf{the question word carries the sense, and the answer carries the case.} \textbf{\sk{कुतः}} itself ends in \textbf{-\sk{तः}}, which is the same \emph{away-from} meaning frozen onto a question word — the question is built out of the answer's grammar.
\end{grammarlens}

\begin{grammarlens}[title={What you've built}]
An answer to a question this book has been asking for forty chapters.
\end{grammarlens}

\subsection*{Your turn}
\begin{itemize}
\raggedright
  \item \emph{Say it:} \emph{kutaḥ?} — then answer it: \emph{Bhāratāt}
  \item \emph{Say it:} \emph{saḥ Bhāratāt āgacchati}
  \item \emph{Say it:} \emph{deśāt}, and say what it means
  \item \emph{Read it:} \textbf{\sk{भारतात्}}, and name the ending that makes it \textquotedblleft{}from\textquotedblright{}
\end{itemize}

\begin{tcolorbox}[breakable,colback=teal!4,colframe=teal!35!black,title={Before you move on}]
Which ending means \textquotedblleft{}from\textquotedblright{}? (\textbf{-\sk{आत्}}.) Answer \textbf{\sk{कुतः}} out loud.
\end{tcolorbox}
